\begin{thebibliography}{99}
\bibitem[Abriola et~al.(2024)]{AbriolaEtAl2024}
Sergio Abriola, Simon Halfon, Aliaume Lopez, Sylvain Schmitz, Philippe
Schnoebelen and Isa Vialard.
\newblock \emph{Measuring well-quasi-ordered finitary powersets}.
\newblock arXiv:2312.14587v2, 17 July 2024.
\newblock Open extension question: Section 6, p.~23.
\newblock \url{https://arxiv.org/abs/2312.14587}.

\bibitem[de Jongh and Parikh(1977)]{deJonghParikh1977}
D.~H.~J. de Jongh and Rohit Parikh.
\newblock Well-partial orderings and hierarchies.
\newblock \emph{Indagationes Mathematicae}, 39(3):195--207, 1977.

\bibitem[D{\v z}amonja et~al.(2020)]{DSS2020}
Mirna D{\v z}amonja, Sylvain Schmitz and Philippe Schnoebelen.
\newblock On ordinal invariants in well quasi orders and finite antichain orders.
\newblock In \emph{Well Quasi-Orders in Computation, Logic, Language and
Reasoning}, Trends in Logic 53, pp.~29--54. Springer, 2020.
\newblock Updated preprint: arXiv:1711.00428v5, 28 January 2025.
\newblock \url{https://arxiv.org/abs/1711.00428}.

\bibitem[Vialard(2025)]{Vialard2025}
Isa Vialard.
\newblock On maximal order type of the lexicographic product.
\newblock \emph{Logic Journal of the IGPL}, 33(4):jzaf051, 2025.
\newblock Published 17 July 2025.
\newblock \url{https://doi.org/10.1093/jigpal/jzaf051}.

% Entries below were added with the batch-36 addition (Parts II--III).
\bibitem[Hilton(2016)]{Hilton2016}
Jacob Hilton.
\newblock The topological pigeonhole principle for ordinals.
\newblock arXiv:1410.2520v4, 12 July 2016.
\newblock Section 2.5, Theorem 2.20 (finite Milner--Rado formula).
\newblock \url{https://arxiv.org/abs/1410.2520}.

\bibitem[Marcone et~al.(2012)]{MMS2012}
Alberto Marcone, Antonio Montalb\'an and Richard~A. Shore.
\newblock Computing maximal chains.
\newblock arXiv:1201.4408, 2012.
\newblock Theorem 1.2 states Wolk's maximal-chain theorem.
\newblock \url{https://arxiv.org/abs/1201.4408}.

\bibitem[Schmidt(1981)]{Schmidt1981}
D.~Schmidt.
\newblock The relation between the height of a well-founded partial ordering
and the order types of its chains and antichains.
\newblock \emph{Journal of Combinatorial Theory, Series B}, 31(2):183--189, 1981.

\bibitem[Vialard(2025b)]{VialardWidth2025}
Isa Vialard.
\newblock On the width of the Cartesian product of ordinals.
\newblock \emph{Order}, 42:37--58, 2025 (published online 2024).
\newblock \url{https://doi.org/10.1007/s11083-024-09668-8}.

\bibitem[Wolk(1967)]{Wolk1967}
E.~S. Wolk.
\newblock Partially well ordered sets and partial ordinals.
\newblock \emph{Fundamenta Mathematicae}, 60(2):175--186, 1967.
\newblock Theorem 9 (maximal chains).
\newblock \url{https://doi.org/10.4064/fm-60-2-175-186}.

% Entries below were added with the batch-42 addition (Part IV).
\bibitem[Bre{\v s}ar et~al.(2016)]{BHR2016}
Bo{\v s}tjan Bre{\v s}ar, Michael~A. Henning and Douglas~F. Rall.
\newblock Total dominating sequences in graphs.
\newblock \emph{Discrete Mathematics}, 339:1665--1676, 2016.
\newblock \url{https://doi.org/10.1016/j.disc.2016.01.017}.
\newblock Author manuscript: \url{https://arxiv.org/abs/1601.07525}.

\bibitem[Francis et~al.(2018)]{FJJ2018}
Mathew~C. Francis, Dalu Jacob and Satyabrata Jana.
\newblock Uniquely restricted matchings in interval graphs.
\newblock \emph{SIAM Journal on Discrete Mathematics}, 32(1):148--172, 2018.
\newblock \url{https://doi.org/10.1137/16M1074631}.
\newblock Author manuscript: \url{https://arxiv.org/abs/1604.07016}.

\bibitem[Golumbic et~al.(2001)]{GHL2001}
Martin~C. Golumbic, Tirza Hirst and Moshe Lewenstein.
\newblock Uniquely restricted matchings.
\newblock \emph{Algorithmica}, 31:139--154, 2001.
\newblock \url{https://doi.org/10.1007/s00453-001-0004-z}.
\end{thebibliography}
